\documentclass[10pt,a4paper]{article}
\usepackage{amsmath,amssymb,amsthm,graphicx}
\usepackage{amsmath,amsfonts,amssymb,amsthm,enumerate,hyperref,ytableau,authblk,listings}
\usepackage[capitalize]{cleveref}
\newtheorem{theorem}{Theorem}[section]
\newtheorem{lemma}[theorem]{Lemma}
\newtheorem{corollary}[theorem]{Corollary}
\newtheorem{conjecture}[theorem]{Conjecture}

\theoremstyle{definition}

\newtheorem{definition}[theorem]{Definition}
\providecommand{\coset}[2]{[#1 \mapsto #2]}
\providecommand{\dcoset}[2]{[#1 \leftrightarrow #2]}
\providecommand{\id}{\mathsf{id}}
\providecommand{\cM}{\mathcal{M}}
\title{A degree-two family of permutations is covered by few cosets of its certificate size}
\author{Gilad Chase, Neta Dafni, Yuval Filmus, Nathan Lindzey}
\date{arxiv:2210.00245v2; Creative Commons Attribution 4.0 International}
\begin{document}
\maketitle
\noindent\textit{Editorial note.} The counted statement is the author's labelled environment at source0.tex lines 608-616 and its proof environment opens at line 617; the pairing is by adjacency and the retained proof range is the whole environment.\par\medskip
\noindent\textit{Editorial note.} Bibliography: the article's own compiled bibliography file (main.bbl, a complete thebibliography with 29 entries) is bundled unchanged; no entry is edited or re-rendered.\par\medskip
\noindent\textit{Editorial note.} Retained-licence form: four-key reuse record in a task-directory evidence copy; no status key and no acceptance record are written.\par\medskip
\noindent\textit{Editorial note.} Metadata provenance: source\_version from the source's licence evidence; authors from the authors\_index key of its source record, title from the title key.\par\medskip
\noindent\textit{Editorial note.} No label-number dump exists for this source and no counter values are imposed, so identity is the author label.\par\medskip
\noindent\textit{Editorial scope.} One original result of the pinned author manuscript, counted once: the coset covering bound for degree-two families of permutations (author label lem:main-ub-sym).\par\medskip
\noindent\textit{Allowed standard prerequisites (editorial).}\begin{itemize}
\item Coset covering. The family is contained in the union of T cosets, each determined by prescribing C(f) images.. It is the conclusion.
\item Representative lemma. Two members with equal value admit representatives meeting exactly in the intersection of their certificates.. It is used to place members in cosets.
\item Structure of a covered family. An m-covered family gives every member a certificate of size m plus one.. It bounds the number of cosets needed.
\end{itemize}\medskip

\begin{lemma} \label{lem:2-int-coset-sym}
Let $f\colon S_n \to \{0,1\}$ be the characteristic function of a $2$-intersecting family, and suppose that $C(f) \leq n-2$. If $f(\alpha) = f(\beta) = 1$ and $C_\alpha,C_\beta$ are minimum certificates for $\alpha,\beta$ (respectively), then $|C_\alpha \cap C_\beta| \geq 2$.
\end{lemma}

\begin{lemma} \label{lem:covering-sym}
Let $f$ be the characteristic function of a $2$-intersecting family $\mathcal{F} \subseteq S_n$, and suppose that $C(f) \leq n-2$.

If $\alpha \in \mathcal{F}$ then $\mathcal{F}$ is $m$-covered for $m = C(f,\alpha) - 1$.
\end{lemma}

\begin{lemma} \label{lem:structure-sym}
Let $\mathcal{F} \subseteq S_n$ be an $m$-covered family, say $\mathcal{F} \subseteq \coset{i_1}{j_1} \cup \cdots \cup \coset{i_m}{j_m}$. Suppose that its characteristic function $f$ has degree at most~$2$.

Every $\alpha$ has a certificate of the form $\{(i_1,\alpha(i_1)),\ldots,(i_m,\alpha(i_m)),(i,j)\}$, and in particular, $C(f) \leq m + 1$.
\end{lemma}

\begin{lemma} \label{lem:main-ub-sym}
Let $\mathcal{F} \subseteq S_n$ be a $2$-intersecting family whose characteristic function $f$ has degree at most~$2$. Suppose that $C(f) \leq n-2$.

The family $\mathcal{F}$ is contained in the union of at most $T$ many $C(f)$-cosets, where
\[
 T = \frac{2\lfloor C(f)/2 \rfloor (C(f) - 1)!}{2^{\lfloor C(f)/2 \rfloor}}.
\]
Moreover, $C(f) \geq 2$.
\end{lemma}

\begin{proof}
If $\mathcal{F}$ is empty then the lemma is vacuously true. Otherwise, choose an arbitrary $\beta \in \mathcal{F}$. \Cref{lem:2-int-coset-sym}, applied to $\alpha = \beta$, shows that $C(f,\beta) \geq 2$. \Cref{lem:covering-sym} shows that $\mathcal{F}$ is $m$-covered, where $m = C(f,\beta) - 1$. Without loss of generality, we can assume that $\mathcal{F} \subseteq \coset{1}{1} \cup \cdots \cup \coset{m}{m}$.

\Cref{lem:structure-sym} shows that $C(f) \leq m + 1$, and so $m = C(f) - 1$ and $C(f) = C(f,\beta) \geq 2$. The \lcnamecref{lem:structure-sym} also shows that every $\alpha \in \mathcal{F}$ has a certificate $C_\alpha$ of the form $\{(1,\alpha(1)),\ldots,(m,\alpha(m)),(i,j)\}$.
It also implies that $\mathcal{F}$ is not $(m-1)$-covered (otherwise its certificate complexity would be at most $m < C(f)$). Therefore, if $X$ is any set of at most $m-1$ compatible cosets, then there is $\alpha_X \in \mathcal{F}$ which is not contained in the union of the cosets in $X$. Fix a certificate $C_X$ for $\alpha_X$ of the form $\{(1,\alpha_X(1)),\ldots,(m,\alpha_X(m)),(i_X,j_X)\}$.

We think of the certificates $C_\alpha,C_X$ as colored using $m + 1$ colors: for $k \in [m]$, a pair $(k,\ast)$ is colored $k$, and the remaining pair $(i,j)$ is colored $m+1$. We can determine the color of a pair $(i',j')$ given only the pair: if $i' \in [m]$ then the color is $i'$, and otherwise the color is $m+1$.

We will prove the lemma by showing that if $\alpha \in \mathcal{F}$ then $C_\alpha$ is one of at most $T$ possible certificates. Since $C_\alpha$ corresponds to a $C(f)$-coset containing $\alpha$, the lemma would follow.
In the proof, we will freely identify a pair $(i,j)$ with the corresponding coset $\coset{i}{j}$, and from now on we will only ever mention cosets.

\smallskip

The proof is slightly different depending on the parity of $C(f)$. We start with the case in which $C(f) = 2r$ is even. Let $\alpha \in \mathcal{F}$.
We define a sequence $X_0 \subset \cdots \subset X_r \subseteq C_\alpha$, where $X_s$ is a collection of $2s$ compatible cosets, as follows. The starting point is $X_0 = \emptyset$. Now let $s < r$, and suppose that $X_s$ has been defined.  Since $|X_s| = 2s \leq C(f) - 2 = m - 1$, there is a permutation $\alpha_{X_s} \in \mathcal{F}$ which lies outside of the union of the cosets in $X_s$. According to \Cref{lem:2-int-coset-sym}, the certificates $C_\alpha$ and $C_{X_s}$ must have at least two cosets in common $q_1,q_2$. The choice of $\alpha_{X_s}$ guarantees that $q_1,q_2 \notin X_s$, and we define $X_{s+1} = X_s \cup \{q_1,q_2\}$. By construction, $X_{s+1} \subseteq C_\alpha$, and so the cosets in $X_{s+1}$ are indeed compatible.

Note that $|X_r| = 2r = |C_\alpha|$, and so $X_r = C_\alpha$. For $s < r$, given $X_s$, the set $X_{s+1} \setminus X_s$ consists of two cosets belonging to $C_{X_s}$. Moreover, since $X_{s+1} \subseteq C_\alpha$ and $C_\alpha$ contains one coset of each color, the colors of the two cosets in $X_{s+1} \setminus X_s$ are different from the colors of the cosets in $X_s$. Since $C_{X_s}$ also contains one coset of each color, this means that there are $\binom{|C_{X_s}| - |X_s|}{2} = \binom{C(f) - 2s}{2}$ choices for $X_{s+1} \setminus X_s$. In total, the number of possible choices for $C_\alpha$ is
\[
 \prod_{s=0}^{r-1} \binom{C(f) - 2s}{2} =
 \frac{C(f)!}{2^r},
\]
matching the formula for $T$ in the statement of the lemma.

\smallskip

Now suppose that $C(f) = 2r+1$ is odd. Let $\alpha \in \mathcal{F}$. We define a sequence $X_0 \subset \cdots \subset X_r \subseteq C_\alpha$, where $X_s$ is a collection of $2s+1$ compatible cosets, as follows. We start by defining $X_0$. According to \Cref{lem:2-int-coset-sym}, the certificates $C_\alpha$ and $C_\emptyset$ must contain at least two cosets in common $q_1,q_2$. At least one of these cosets is of the form $(k,\alpha(k))$ for some $k \in [m]$. We define $X_0 = \{(k,\alpha(k))\}$. Now let $s < r$, and suppose that $X_s$ has been defined. Since $|X_s| = 2s+1 \leq C(f) - 2 = m - 1$, we can define $X_{s+1}$ as in the case of even $C(f)$.

Note that $|X_r| = 2r+1 = |C_\alpha|$, and so $X_r = C_\alpha$. The certificate $X_0$ contains a single coset $(k,\ell) \in C_\emptyset$, where $k \in [m]$, and so there are $m = C(f) - 1$ choices for $X_0$. For $s < r$, given $X_s$, there are $\binom{|C_{X_s}| - |X_s|}{2} = \binom{C(f) - 2s - 1}{2}$ choices for $X_{s+1}$ (this is the same as in the case of even $C(f)$). In total, the number of possible choices for $C_\alpha$ is
\[
 (C(f) - 1) \prod_{s=0}^{r-1} \binom{C(f) - 2s - 1}{2} =
 (C(f) - 1) \frac{(C(f) - 1)!}{2^r},
\]
matching the formula for $T$ in the statement of the lemma.
\end{proof}

\newcommand{\etalchar}[1]{$^{#1}$}
\begin{thebibliography}{BMRR21}

\bibitem[AK97]{AK1}
Rudolf Ahlswede and Levon~H. Khachatrian.
\newblock The complete intersection theorem for systems of finite sets.
\newblock {\em European J. Combin.}, 18(2):125--136, 1997.

\bibitem[AK99]{AK2}
Rudolf Ahlswede and Levon~H. Khachatrian.
\newblock A pushing-pulling method: new proofs of intersection theorems.
\newblock {\em Combinatorica}, 19(1):1--15, 1999.

\bibitem[BdW02]{BuhrmandeWolf}
Harry Buhrman and Ronald de~Wolf.
\newblock Complexity measures and decision tree complexity: a survey.
\newblock {\em Theoret. Comput. Sci.}, 288(1):21--43, 2002.
\newblock Complexity and logic (Vienna, 1998).

\bibitem[BMRR21]{BMRR}
Angelot Behajaina, Roghayeh Maleki, Aina~Toky Rasoamanana, and A.~Sarobidy
  Razafimahatratra.
\newblock 3-setwise intersecting families of the symmetric group.
\newblock {\em Discrete Math.}, 344(8):Paper No. 112467, 15, 2021.

\bibitem[CK03]{CameronKu}
Peter~J. Cameron and C.~Y. Ku.
\newblock Intersecting families of permutations.
\newblock {\em European J. Combin.}, 24(7):881--890, 2003.

\bibitem[CSST08]{CSST08}
Tullio Ceccherini-Silberstein, Fabio Scarabotti, and Filippo Tolli.
\newblock {\em Harmonic analysis on finite groups}, volume 108 of {\em
  Cambridge Studies in Advanced Mathematics}.
\newblock Cambridge University Press, Cambridge, 2008.
\newblock Representation theory, Gelfand pairs and Markov chains.

\bibitem[Daf22]{DafniThesis}
Neta Dafni.
\newblock Complexity measures on the symmetric group and beyond.
\newblock Master's thesis, Technion --- Israel Institute of Technology, 2022.

\bibitem[DFL{\etalchar{+}}21]{DFLLV}
Neta Dafni, Yuval Filmus, Noam Lifshitz, Nathan Lindzey, and Marc Vinyals.
\newblock Complexity measures on the symmetric group and beyond.
\newblock In {\em 12th {I}nnovations in {T}heoretical {C}omputer {S}cience
  {C}onference}, volume 185 of {\em LIPIcs. Leibniz Int. Proc. Inform.}, pages
  Art. No. 87, 5. Schloss Dagstuhl. Leibniz-Zent. Inform., Wadern, 2021.

\bibitem[EFP11]{EFP}
David Ellis, Ehud Friedgut, and Haran Pilpel.
\newblock Intersecting families of permutations.
\newblock {\em J. Amer. Math. Soc.}, 24(3):649--682, 2011.

\bibitem[EKR61]{EKR}
P.~Erd\H{o}s, Chao Ko, and R.~Rado.
\newblock Intersection theorems for systems of finite sets.
\newblock {\em Quart. J. Math. Oxford Ser. (2)}, 12:313--320, 1961.

\bibitem[Ell12a]{EllisCameronKu}
David Ellis.
\newblock A proof of the {C}ameron-{K}u conjecture.
\newblock {\em J. Lond. Math. Soc. (2)}, 85(1):165--190, 2012.

\bibitem[Ell12b]{EllisSetwise}
David Ellis.
\newblock Setwise intersecting families of permutations.
\newblock {\em J. Combin. Theory Ser. A}, 119(4):825--849, 2012.

\bibitem[FD77]{DezaFrankl}
P\'{e}ter Frankl and Mikhail Deza.
\newblock On the maximum number of permutations with given maximal or minimal
  distance.
\newblock {\em J. Combinatorial Theory Ser. A}, 22(3):352--360, 1977.

\bibitem[FHKL16]{FHKL}
Yuval Filmus, Hamed Hatami, Nathan Keller, and Noam Lifshitz.
\newblock On the sum of the {$L_1$} influences of bounded functions.
\newblock {\em Israel J. Math.}, 214(1):167--192, 2016.

\bibitem[Fil17]{FilmusComment}
Yuval Filmus.
\newblock A comment on intersecting families of permutations.
\newblock arXiv:1706.10146, 2017.

\bibitem[FMS21]{FMS}
Shaun Fallat, Karen Meagher, and Mahsa~N. Shirazi.
\newblock The {E}rd{\H{o}}s-{K}o-{R}ado theorem for 2-intersecting families of
  perfect matchings.
\newblock {\em Algebr. Comb.}, 4(4):575--598, 2021.

\bibitem[GM09]{GodsilMeagher}
Chris Godsil and Karen Meagher.
\newblock A new proof of the {E}rd{\H{o}}s-{K}o-{R}ado theorem for intersecting
  families of permutations.
\newblock {\em European J. Combin.}, 30(2):404--414, 2009.

\bibitem[GM16]{GodsilMeagherBook}
Chris Godsil and Karen Meagher.
\newblock {\em Erd\H{o}s-{K}o-{R}ado theorems: algebraic approaches}, volume
  149 of {\em Cambridge Studies in Advanced Mathematics}.
\newblock Cambridge University Press, Cambridge, 2016.

\bibitem[GM17]{GodsilMeagherPM}
Chris Godsil and Karen Meagher.
\newblock An algebraic proof of the {E}rd{\H{o}}s-{K}o-{R}ado theorem for
  intersecting families of perfect matchings.
\newblock {\em Ars Math. Contemp.}, 12(2):205--217, 2017.

\bibitem[Lin17]{Lindzey17}
Nathan Lindzey.
\newblock Erd{\H{o}}s-{K}o-{R}ado for perfect matchings.
\newblock {\em European J. Combin.}, 65:130--142, 2017.

\bibitem[Lin18a]{Lindzey18}
Nathan Lindzey.
\newblock Intersecting families of permutations of perfect matchints.
\newblock arXiv:1811.06160, 2018.

\bibitem[Lin18b]{LindzeyThesis}
Nathan Lindzey.
\newblock {\em Matchings and Representation Theory}.
\newblock PhD thesis, University of Waterloo, 2018.

\bibitem[LM04]{LaroseMalvenuto}
Benoit Larose and Claudia Malvenuto.
\newblock Stable sets of maximal size in {K}neser-type graphs.
\newblock {\em European J. Combin.}, 25(5):657--673, 2004.

\bibitem[MM05]{MeagherMoura}
Karen Meagher and Lucia Moura.
\newblock Erd{\H{o}}s-{K}o-{R}ado theorems for uniform set-partition systems.
\newblock {\em Electron. J. Combin.}, 12:Research Paper 40, 12, 2005.

\bibitem[MR21]{MR21}
Karen Meagher and A.~Sarobidy Razafimahatratra.
\newblock The {E}rd{\H{o}}s-{K}o-{R}ado theorem for 2-pointwise and 2-setwise
  intersecting permutations.
\newblock {\em Electron. J. Combin.}, 28(4):Paper No. 4.10, 21, 2021.

\bibitem[N{\"O}03]{Cliquer}
Sampo Niskanen and Patric R.~J. {\"O}sterg{\aa}rd.
\newblock Cliquer user's guide, version 1.0.
\newblock Technical Report T48, Communications Laboratory, Helsinki University
  of Technology, Espoo, Finland, 2003.

\bibitem[O'D14]{ODonnell}
Ryan O'Donnell.
\newblock {\em Analysis of {B}oolean functions}.
\newblock Cambridge University Press, New York, 2014.

\bibitem[Sag01]{Sagan}
Bruce~E. Sagan.
\newblock {\em The symmetric group}, volume 203 of {\em Graduate Texts in
  Mathematics}.
\newblock Springer-Verlag, New York, second edition, 2001.
\newblock Representations, combinatorial algorithms, and symmetric functions.

\bibitem[{Sag}22]{sagemath}
{Sage Developers}.
\newblock {\em {S}ageMath, the {S}age {M}athematics {S}oftware {S}ystem
  ({V}ersion 9.5)}, 2022.
\newblock {\tt https://www.sagemath.org}.

\end{thebibliography}
\end{document}
